\section{Background you'll need}
\label{sec:background}

This report is written to be read by someone with a first course in statistics, a working idea of
what fine-tuning a small language model means, and a lay familiarity with mortgage lending. Everything
technical the four acts rely on is defined once here, from the ground up, with the exact equations and
a worked example wherever a number would otherwise be mysterious. Nothing in this section reports a
result; it is the vocabulary.

\begin{background}{The one-paragraph orientation}
A \emph{guard} is a small model placed in front of an assistant: it reads an incoming request and
labels it \code{safe} or \code{unsafe} before anything acts on it. We build each guard by
\emph{fine-tuning} a general chat model, and we score it by how well its numeric output \emph{ranks}
truly-unsafe requests above safe ones. The central move of this report is to always compare a
fine-tuned guard against the \emph{same model before tuning}, on identical inputs, and to separate two
kinds of test set: benchmarks whose data the guard \emph{trained on} (represented) and benchmarks it
\emph{never saw} (transfer). Almost every subtlety below is a way of making that comparison honest.
\end{background}

\subsection{The guard: a single-token logit-difference head}
\label{sec:bg-guard}

We do not ask the guard to write a paragraph explaining its decision. We ask it for exactly one
word and look only at how strongly it leans toward that word. Concretely, a prompt $x$ is wrapped in a
fixed instruction template asking for a one-word verdict, and we read the model's raw output scores
--- its \emph{logits} --- at the final position for the two verdict words \code{safe} and
\code{unsafe}. A logit is the model's un-normalized, pre-probability score for a candidate next token:
larger means the model favors that token. Writing $z_{\text{unsafe}}(x,t)$ and $z_{\text{safe}}(x,t)$
for those two logits at the final prompt position $t$, the guard's stored score is their difference,
\begin{equation}
\label{eq:score}
s(x) \;=\; z_{\text{unsafe}}(x,t) \;-\; z_{\text{safe}}(x,t),
\end{equation}
which is positive when the model leans \code{unsafe}, negative when it leans \code{safe}, and near zero
when it is undecided. To read $s(x)$ as a probability we pass the same two logits through a two-way
softmax,
\begin{equation}
\label{eq:prob}
p(\text{unsafe}\mid x) \;=\; \frac{\exp\!\big(z_{\text{unsafe}}\big)}{\exp\!\big(z_{\text{safe}}\big)+\exp\!\big(z_{\text{unsafe}}\big)}.
\end{equation}
For example, if $z_{\text{unsafe}}=2.0$ and $z_{\text{safe}}=1.0$, then $s(x)=1.0$ and
$p(\text{unsafe}\mid x)=e^{2}/(e^{2}+e^{1})\approx 0.73$. This is what we mean by a
``single-token logit-difference head'': one forward pass, two numbers, one score. The raw logits are
the canonical stored value; probabilities in \Cref{eq:prob} and any temperature-rescaled version are
\emph{derived} from them. Reading the model's own generated verdict text is deliberately never used as
the primary score, because a free-text answer is harder to score consistently and hides the model's
actual confidence.

\subsection{Base, SFT, and LoRA}
\label{sec:bg-sft}

The \emph{base} is the general instruction-tuned chat model \emph{before} we turn it into a guard. It
is not untrained --- it can already follow the ``answer \code{safe} or \code{unsafe}'' instruction
zero-shot --- so ``base'' throughout means ``before the guard fine-tune,'' not ``blank.''
\emph{Supervised fine-tuning} (SFT) then trains that model on labeled examples: for each training
prompt we know the correct verdict, and we nudge the model to raise the probability it assigns to that
verdict word. Rather than update all of the model's weights, we use \emph{LoRA} (Low-Rank Adaptation)
\citep{hu2021lora}: we freeze the original weights and insert a small number of extra trainable
parameters (a low-rank ``adapter'') into selected weight matrices. LoRA is cheap, fast, and reversible
--- you can ship the tiny adapter on top of the frozen base --- and it is the standard way guards are
built in practice \citep{elesedy2024loraguard}. The precise adapter size and which matrices it touches
are part of the frozen recipe in \Cref{sec:setup}.

\subsection{Average precision, macro-AP, and the accuracy trap}
\label{sec:bg-ap}

Because $s(x)$ is a continuous score, we can grade a guard \emph{without} committing to a cutoff, by
asking how well it \emph{ranks}. \emph{Average precision} (AP) collapses the whole
precision--recall curve into a single number in $[0,1]$: it is near $1$ when the guard tends to give
genuinely unsafe prompts higher scores than safe ones, and it needs no threshold. Formally, sweeping
the cutoff from high to low, AP is the recall-weighted average of the precision achieved at each
true-positive, $\text{AP}=\sum_n (R_n-R_{n-1})\,P_n$, with $P_n,R_n$ the precision and recall after the
$n$-th ranked item; we use the tie-aware, non-interpolated \code{scikit-learn} definition so the number
is exactly reproducible.

\begin{background}{A worked AP example}
Rank five prompts by their guard score, highest first, and write their true labels:
$[\text{unsafe},\ \text{unsafe},\ \text{safe},\ \text{unsafe},\ \text{safe}]$. There are three unsafe
prompts, so each one contributes recall $\tfrac13$. The precision at each unsafe prompt as we go down
the list is $\tfrac11{=}1.00$ (rank 1), $\tfrac22{=}1.00$ (rank 2), and $\tfrac34{=}0.75$ (rank 4). So
$\text{AP}=\tfrac13(1.00)+\tfrac13(1.00)+\tfrac13(0.75)\approx 0.917$. A perfect ranking (all three
unsafe prompts on top) would score $1.00$; ranking the safe prompts first would score much lower.
\end{background}

\emph{Macro}-AP computes AP separately on each benchmark and then averages those per-benchmark values
with \emph{equal weight}, so one large benchmark cannot dominate a handful of small ones; this is the
primary metric in every act.

\begin{takeaway}{Why not just report accuracy?}
On a test set that is $95\%$ safe, a guard that blindly answers \code{safe} to everything scores $95\%$
accuracy while catching \emph{zero} attacks. Accuracy rewards guessing the majority class; AP and
macro-AP instead reward putting the unsafe prompts on top, which is the behavior a guard is for. This
is why ranking, not accuracy, is the currency of this report.
\end{takeaway}

\subsection{Two evaluation regimes: represented-source vs.\ dataset-held-out transfer}
\label{sec:bg-regimes}

The single most important distinction in this report is \emph{where the test data came from}.
\textbf{Represented-source} test rows are held-back examples drawn from a dataset the guard
\emph{trained on} (different rows, same source). \textbf{Dataset-held-out transfer} rows come from
datasets whose examples were \emph{never used in training at all}. A gain that appears on
represented-source tests but not on transfer tests is precisely the signature of a guard
\emph{specializing} to its training sources rather than becoming a broadly better moderator. One
honesty caveat travels with the word ``held out'': it means \emph{dataset-held-out} (those rows were
withheld from fine-tuning), \emph{not} sealed away from the researchers during development --- a
distinction that forces the retrospective framing discussed in \Cref{sec:bg-estimand}.

\subsection{From scores to decisions: calibration, operating point, TPR/FPR, macro vs.\ pooled}
\label{sec:bg-op}

Ranking is threshold-free, but a deployed guard must eventually say \code{safe} or \code{unsafe}. An
\emph{operating point} is the single cutoff on $s(x)$ that converts scores into hard decisions. Picking
it trades two error rates against each other: the \emph{true-positive rate} (TPR, or recall) is the
fraction of genuinely unsafe prompts the guard flags, and the \emph{false-positive rate} (FPR) is the
fraction of genuinely safe prompts it wrongly flags. Raising the cutoff lowers both; lowering it raises
both. Before thresholding we \emph{calibrate}: we fit a single positive \emph{temperature} that rescales
the logits (temperature scaling \citep{guo2017calibration}) so the derived probabilities in
\Cref{eq:prob} are better behaved. Crucially, both the temperature and the cutoff are fit only on a
separate held-back \emph{calibration} split --- never on the transfer or stress data we report on --- so
the threshold is not quietly tuned to the test.

Two error rates can be averaged two ways, and the choice matters. \emph{Macro} rates compute the rate
on each benchmark first and then average across benchmarks (equal weight per benchmark). \emph{Pooled}
rates lump every row together first and then compute one rate (equal weight per row, so large
benchmarks dominate). We report both because they can disagree.

\begin{takeaway}{Ranking $\neq$ a transferable threshold}
A guard can rank well (high AP) yet have no cutoff that holds up off-source: the score
\emph{distribution} shifts even when the \emph{ordering} is fine. So AP answers ``does it rank?'' and
the operating point answers ``is there a usable decision boundary?''; the two questions have different
answers, and a deployment claim from AP alone would miss the gap.
\end{takeaway}

\subsection{Estimands, the fixed panel, and the paired hierarchical bootstrap}
\label{sec:bg-estimand}

An \emph{estimand} is the exact quantity our numbers describe. Here it is the average
before-versus-after change \emph{for the four specific checkpoints we study}, not for ``small guards''
in general. Those four checkpoints are a \emph{fixed, purposively chosen panel} --- picked deliberately,
not sampled at random from a population --- so we do not attach uncertainty to the choice of models and
we do not generalize to unnamed architectures. Every comparison is \emph{paired}: an adapter is always
compared against \emph{its own base} on the \emph{same rows}, so nothing but the fine-tune differs.

To put a $95\%$ range around a paired change we use a \emph{paired hierarchical bootstrap}. A bootstrap
estimates how much a result would wobble under resampling by rebuilding the data many times from itself
and recomputing; the spread of the recomputed values becomes the interval. Ours is \emph{hierarchical}
because it resamples at two levels at once --- which fine-tuning \emph{seeds} were drawn (within each
fixed checkpoint), and which groups of near-duplicate evaluation rows (``families,''
\Cref{sec:setup-decon}) are counted, via one Poisson(1) re-count weight per family. It is \emph{paired}
because each adapter stays yoked to its base. Because the training rows and their order are held fixed,
seed-to-seed variation reflects initialization, dropout, and execution nondeterminism --- not
which examples the guard happened to see. The resulting intervals are \emph{conditional} on this panel,
these datasets, and this family graph; they are descriptive, not significance tests, and not claims
about a wider population.

\begin{takeaway}{Retrospective, estimation-only}
The fine-tuning studies (Acts I--III) were analyzed \emph{after} the benchmarks had been inspected
during development, so we report point estimates, intervals, and leave-one-out sensitivity checks ---
but no accept/reject hypothesis test and no ``pass/fail gate.'' Clean provenance does not make an
inspected benchmark prospective. Confirmatory use would require a freshly locked protocol on a
genuinely uninspected cohort.
\end{takeaway}

\subsection{Three fine-tuning objectives: SFT, DPO, GRPO}
\label{sec:bg-objectives}

Act~II asks whether the \emph{objective} used to fine-tune changes how much a guard specializes. Three
objectives are compared. \textbf{SFT} minimizes cross-entropy: it directly raises the model's log
probability of the single correct verdict token. \textbf{DPO} (Direct Preference Optimization
\citep{rafailov2023dpo}) learns from a \emph{pair} --- the correct verdict versus the wrong one --- and
increases the correct verdict's relative log-probability margin under a reference-model regularizer, so
it never needs a separately trained reward model. \textbf{GRPO} (Group Relative Policy Optimization,
introduced with DeepSeekMath \citep{shao2024deepseekmath}) is reinforcement learning in which the
\emph{label itself is a verifiable reward} (correct verdict earns reward, otherwise not) and advantages
are computed relative to a group of sampled outputs --- again with no learned reward model. A structural
caveat matters for guards and is pre-registered in Act~II: because our action space is a
\emph{single token with two outcomes}, the exploration that lets RL generalize in richer generation
tasks is largely absent, so RL's usual transfer edge may be small or null here
\citep{vennemeyer2026objective}.

\subsection{Output-space composition}
\label{sec:bg-composition}

Act~III introduces a retraining-free remedy. \emph{Output-space composition} combines two guards by
averaging their calibrated \emph{output scores} --- not their internal weights. It is portable, because
it needs only two comparable scores rather than interpolable parameters, but it costs two inference
passes because both models must run. This contrasts with \emph{weight-space} methods such as WiSE-FT and
model soups \citep{wortsman2022wiseft,wortsman2022modelsoups}, which average the models' weights and so
require the models to be weight-compatible. The exact composition rule is given later as
\Cref{eq:comp}.

\subsection{The mortgage/HMDA vocabulary and the dual-label idea}
\label{sec:bg-mortgage}

Act~IV moves to regulated domains, where a request can read as perfectly polite text yet, if honored,
break the law. The mortgage vocabulary below is what the benchmark encodes.

\begin{background}{Mortgage terms, in plain words}
\textbf{HMDA} --- the U.S.\ Home Mortgage Disclosure Act \citep{hmda2022}: lenders publicly report
loan-level records (purpose, amount band, action taken, denial reason), which we use only as
\emph{de-identified, banded fact sheets}, never verbatim. \textbf{Fair lending} bars credit
discrimination on protected traits. \textbf{Redlining} is refusing or worsening credit across a whole
neighborhood as a proxy for race. \textbf{Disparate treatment} is intentional differential treatment;
\textbf{disparate impact} is a facially neutral rule that falls harder on a protected group.
\textbf{Proxy (coded) discrimination} uses a stand-in --- ZIP code, preferred language --- for a
protected class. \textbf{ATR/QM} is the ability-to-repay / qualified-mortgage rule.
\textbf{Adverse-action notice} is the legally required true reason for a denial. \textbf{UDAAP} covers
unfair, deceptive, or abusive acts and practices.\\[2pt]
\textbf{The dual-label construct.} Each request carries two \emph{independent} labels:
$\mathbf{G}$ --- would an ordinary general-safety guard call this \code{unsafe}? --- and $\mathbf{D}$
--- would honoring it break mortgage policy? The load-bearing stratum is \textbf{G0/D1}: it reads
\emph{safe} to a general guard yet is a compliance \emph{violation}. That is exactly the case a
general-safety score cannot see.
\end{background}

\section{The shared experimental setup}
\label{sec:setup}

Every study in this report reuses one fixed, purposively chosen panel of four checkpoints and one
frozen, decontaminated $1{,}200$-row training manifest, scored by the identical single-token
$z_{\text{unsafe}}-z_{\text{safe}}$ head (\Cref{eq:score}) and the same tie-aware macro-AP, under
fail-closed provenance \emph{locks} that bind the data, code, and software versions and refuse to run
if any fingerprint mismatches. This single-panel, single-manifest discipline is what makes SFT, DPO,
GRPO, composition, and the domain evaluations \emph{directly comparable} rather than three stapled
results. The \emph{same} four checkpoints recur across all four acts, and \textbf{Qwen3-4B is the
recurring character}: the strongest base, and --- as the acts will show --- the one that specializes
most, the one composition helps least, and the best-and-fairest zero-shot mortgage guard.

\subsection{The fixed four-checkpoint panel and pinned revisions}
\label{sec:setup-panel}

The panel spans two model lineages and $1.5$--$4$B parameters: Qwen2.5-1.5B-Instruct,
SmolLM2-1.7B-Instruct, SmolLM3-3B \citep{bakouch2025smollm3}, and Qwen3-4B. Each checkpoint is pinned to
a fixed upstream revision, and for each we verify that \code{safe} and \code{unsafe} map to
\emph{distinct single tokens} under a fixed convention (leading-space \code{" safe"}/\code{" unsafe"}
first, no-leading-space fallback), recording the token IDs and strings. \Cref{tab:panel-recipe} lists
the panel and the frozen recipe together. The estimand is this exact four-checkpoint panel; results are
not extrapolated to other scales, architectures, or vendors.

\begin{table}[htbp]
\centering\small
\caption{The fixed model panel and the frozen clean-v2 LoRA-SFT recipe (shared by Acts~I--III).
Revisions are pinned; decision tokens are verified as distinct single tokens per checkpoint.}
\label{tab:panel-recipe}
\resizebox{\textwidth}{!}{%
\begin{tabular}{lll l}
\toprule
Checkpoint & Params & Revision (pinned, abbreviated) & Decision tokens \\
\midrule
\code{Qwen/Qwen2.5-1.5B-Instruct}        & 1.5B & \code{989aa79\ldots71aa306}  & \{\,safe, unsafe\,\} single-token \\
\code{HuggingFaceTB/SmolLM2-1.7B-Instruct} & 1.7B & \code{31b70e2\ldots46d38674} & \{\,safe, unsafe\,\} single-token \\
\code{HuggingFaceTB/SmolLM3-3B}          & 3B   & \code{a07cc9a\ldots335b0ac1} & \{\,safe, unsafe\,\} single-token \\
\code{Qwen/Qwen3-4B}                     & 4B   & \code{1cfa9a7\ldots03b3df60c} & \{\,safe, unsafe\,\} single-token \\
\midrule
\multicolumn{4}{l}{\emph{Recipe (identical for all four bases):} completion-only SFT loss on the verdict $+$ appended EOS;}\\
\multicolumn{4}{l}{LoRA $r{=}32$, $\alpha{=}64$, dropout $0.05$ on \{q,k,v,o,gate,up,down\}; $300$ optimizer steps;}\\
\multicolumn{4}{l}{lr $2\times10^{-4}$ cosine schedule, warmup $0.03$; effective batch $4$; max sequence length $1{,}024$;}\\
\multicolumn{4}{l}{training seeds $42$--$46$ (five per checkpoint), shared data-order seed $42$.}\\
\bottomrule
\end{tabular}}
\end{table}

\subsection{The single-token guard formulation and prompt rendering}
\label{sec:setup-guard}

Every base and every adapter is scored by the same guard formulation from \Cref{sec:bg-guard}: the
prompt is wrapped in one versioned instruction template asking for a one-word verdict, and the last
position's \code{safe}/\code{unsafe} logits are read and stored as $s(x)$ (\Cref{eq:score}), with the
softmax probability (\Cref{eq:prob}) derived from them. The \emph{semantic} system instruction is
shared verbatim across all four checkpoints, but each model's own chat template renders it differently,
producing three distinct rendered-template fingerprints across the panel; each is recorded. Long user
content is budgeted and truncated \emph{before} final rendering, and the runner then asserts that the
complete classification instruction and wrapper survive the truncation --- a contract check motivated by
an earlier artifact that could left-truncate the instruction itself. Truncation strategy and scored
token counts are stored per row.

\subsection{The frozen LoRA-SFT recipe}
\label{sec:setup-recipe}

The recipe (lower panel of \Cref{tab:panel-recipe}) is fixed identically across all four bases so that
the only thing varying within Act~I is the checkpoint. It uses a \emph{completion-only} loss --- the
cross-entropy is applied only to the verdict token and an appended end-of-sequence marker, not to the
prompt --- with a LoRA adapter of rank $r{=}32$ and scaling $\alpha{=}64$ (dropout $0.05$) inserted into
the attention and MLP projections (q, k, v, o, gate, up, down). Training runs $300$ optimizer steps at
learning rate $2\times10^{-4}$ on a cosine schedule with $0.03$ warmup and an effective batch size of
$4$, at maximum sequence length $1{,}024$. At effective batch $4$, the $1{,}200$-row manifest yields
exactly $300$ updates --- one complete exposure of the data. Each base is adapted with five training
seeds ($42$--$46$) that share one data-order seed ($42$), giving $20$ adapters; the four untuned bases
are scored once and reused, for $24$ model bundles in all.

\subsection{The 1,200-row training manifest and exclusions}
\label{sec:setup-manifest}

Training reads one frozen manifest and never resamples a live dataset mid-run. It draws $400$ rows each
from three represented sources --- ToxicChat \citep{lin2023toxicchat}, Prompt-Injections, and
Jailbreak-Classification --- with $200$ safe and $200$ unsafe rows per source, selected by a hashed rank
salted with the data seed and frozen as one shared row order across every checkpoint and seed
(\Cref{tab:manifest}). Two datasets are \emph{deliberately excluded from training}: BeaverTails
\citep{ji2023beavertails}, because its safety annotation labels the prompt--response interaction rather
than the prompt alone, and OR-Bench \citep{cui2024orbench}, which is reserved for the benign stress set;
including either would confound the prompt-only, dataset-held-out design. The study uses the
non-commercial academic data branch (ToxicChat is CC BY-NC 4.0), so any released adapter inherits a
non-commercial mark, and no third-party raw text is redistributed --- only pinned identifiers,
revisions, and content hashes.

\begin{table}[htbp]
\centering\small
\caption{The frozen $1{,}200$-row training manifest (Acts~I--III). All training rows come from three
represented sources; each contributes $400$ rows balanced $200{:}200$. BeaverTails and OR-Bench are
excluded from training by design.}
\label{tab:manifest}
\begin{tabular}{l l r}
\toprule
Represented source & Native label origin & Train rows (safe : unsafe) \\
\midrule
ToxicChat (\code{lmsys/toxic-chat})           & prompt toxicity   & $400$ ($200{:}200$) \\
Prompt-Injections (\code{deepset/prompt-injections}) & prompt injection & $400$ ($200{:}200$) \\
Jailbreak-Classification (\code{jackhhao/\ldots}) & prompt jailbreak & $400$ ($200{:}200$) \\
\midrule
\textbf{Total} & & $\mathbf{1{,}200}$ ($600{:}600$) \\
\bottomrule
\end{tabular}
\end{table}

\subsection{Three evaluation regimes and decontamination}
\label{sec:setup-decon}

Evaluation is partitioned into three regimes, forming a spectrum of how ``new'' each test is to the
guard.
\begin{itemize}[leftmargin=1.4em,itemsep=1pt]
  \item \textbf{Represented-source} $=$ \{ToxicChat, Prompt-Injections, Jailbreak-Classification\} rows
  \emph{held back} from training. These measure discrimination on sources the guard trained on.
  \item \textbf{Dataset-held-out transfer} $=$ \{JailbreakBench \citep{chao2024jailbreakbench}, XSTest
  \citep{rottger2024xstest}, WildGuardTest \citep{han2024wildguard}, WildJailbreak
  \citep{jiang2024wildteaming}\}, whose rows were withheld from SFT. Because these encode heterogeneous
  native constructs (harm, refusal, jailbreak, contrast), their macro-average mixes distinct policies;
  we study \emph{dataset-held-out benchmark transfer}, not one fixed universal policy.
  \item \textbf{Stress} $=$ \{OR-Bench-Hard benign portion (false-positive rate only), HarmBench
  \citep{mazeika2024harmbench} (recall only)\}. These are single-class probes: \textbf{no AP or AUROC is
  computed on stress data}; they only watch for over-blocking (benign FPR) and missed attacks (recall).
\end{itemize}
The primary metric within each regime is benchmark-macro AP; the secondary metric is the
calibration-targeted operating point of \Cref{sec:setup-stats}.

\paragraph{Decontamination (family-isolated splits, MinHash).} To keep near-duplicate rows from
straddling the train/evaluation and calibration/test boundaries, the builder first preserves the
authoritative upstream conversation, pair, and scenario identifiers, then adds deterministic character
$5$-gram \emph{MinHash} edges between near-duplicate rows. It forms the connected components of that
graph --- ``families'' --- and assigns \emph{whole families} to a single split, so that no family can
appear in both training and any reported evaluation, or in both calibration and any reported test or
stress surface. The build \emph{fails closed} if any family crosses a forbidden boundary. This family
graph also supplies the Poisson(1) re-count weights used by the bootstrap (\Cref{sec:bg-estimand}). An
independent audit of $24$ hard assertions covers the source exclusions, schema and row identity,
exact and conflicting-label overlap, pinned revisions and content hashes, selection provenance, family
disjointness, licenses, and near-duplicate dispositions.

\subsection{Estimands and the statistical protocol}
\label{sec:setup-stats}

For regime $R$, checkpoint $b$, and seed $r$, the per-regime metric is the equal-weight macro-AP over
that regime's benchmarks,
\begin{equation}
\label{eq:regime-metric}
M_R(b,r) \;=\; \frac{1}{|K_R|}\sum_{k\in K_R}\mathrm{AP}_k(b,r),
\end{equation}
using the tie-aware, non-interpolated \code{scikit-learn} AP. The paired base-to-tuned change for one
cell is
\begin{equation}
\label{eq:delta}
\Delta_R(b,r) \;=\; M_R(\mathrm{SFT}_{b,r}) \;-\; M_R(\mathrm{base}_b),
\end{equation}
and the fixed-panel aggregate averages the four per-checkpoint deltas \emph{without} resampling
checkpoint identities,
\begin{equation}
\label{eq:agg}
\bar\Delta_R \;=\; \frac{1}{4}\sum_{b}\Big[\tfrac{1}{5}\textstyle\sum_r M_R(\mathrm{SFT}_{b,r})
\;-\; M_R(\mathrm{base}_b)\Big].
\end{equation}
The headline object is the joint vector $\theta=(\bar\Delta_{\text{represented}},
\bar\Delta_{\text{transfer}})$, kept as a two-dimensional quantity and read off the
\emph{specialization plane} of \Cref{fig:plane} --- horizontal axis the represented change, vertical the
transfer change, four quadrants (uniform gain, specialize, uniform loss, transfer-favored) --- rather
than collapsed into a single ``specialization score.'' The axes and quadrant interpretation are
result-independent.

\paragraph{Uncertainty.} We attach $95\%$ two-sided intervals with $10{,}000$ paired
hierarchical-bootstrap replicates (fixed RNG seed $20260712$): checkpoint identities are held fixed, SFT
seed indices are resampled \emph{within} each checkpoint, and one Poisson(1) weight is drawn per
recorded family (\Cref{sec:setup-decon}). Leave-one-checkpoint-out and leave-one-transfer-benchmark-out
values are \emph{deterministic} sensitivity checks, not resampling or population inference. The locked
analysis mode is \emph{precision-focused}: we report estimates, intervals, and sensitivities, with no
intersection--union test, multiplicity correction, bootstrap $p$-value, or pass/fail gate.

\paragraph{Calibration-only operating point.} As a secondary deployment diagnostic, we fit one positive
temperature per bundle on the calibration split, then choose the threshold that maximizes calibration
recall subject to a one-sided $95\%$ Clopper--Pearson \emph{upper} bound of at most $5\%$ on pooled
calibration-negative FPR; if no cutoff qualifies, the cell reports \code{NO\_FEASIBLE\_THRESHOLD} as an
outcome rather than relaxing the target. Temperature and threshold are fit \emph{only} on calibration
rows --- never on transfer or stress data --- and the audit hard-asserts that calibration shares no
family with any reported test or stress surface. The Clopper--Pearson bound assumes independent
Bernoulli negatives, so it is a pooled diagnostic, not a family-aware production guarantee. The
resulting operating-point rates appear in \Cref{tab:sensitivity} (Act~I) and \Cref{tab:pilot-operating}
(Act~III).

\subsection{The objective axis (SFT vs.\ DPO vs.\ GRPO): design}
\label{sec:setup-objective}

Act~II reuses \emph{everything} above --- the same four-checkpoint panel, the same $1{,}200$-row
manifest, the same seeds, the same single-token scorer, and the same macro-AP plus bootstrap machinery
--- and varies \emph{only the fine-tuning objective}, so any difference is attributable to the objective
rather than to data, capacity, or evaluation. SFT is the Act~I arm; DPO and GRPO are the two additional
arms, both trained against the \emph{verifiable label} with no learned reward model
(\Cref{sec:bg-objectives}). The hypotheses and decision rules are \textbf{pre-registered and bound into
the run's lock before execution}: \textbf{H1}, objectives differ mainly on \emph{transfer}, ordered by
how far each moves the model from its base; \textbf{the GRPO single-token null}, that with a
one-token/two-outcome action space RL's transfer edge over SFT is small or null (a bounded negative
result is the \emph{expected, reportable} outcome, not a failure); and \textbf{H2}, that composition
(\Cref{sec:bg-composition}) recovers transfer in proportion to how far the objective moved from base.
\textbf{Results are \textsc{pending} the locked GPU run and are reported in \Cref{sec:actII}; no
objective-axis numbers appear until its per-row scores are committed.}

\subsection{The composition protocol}
\label{sec:setup-composition}

Act~III's remedy needs no retraining. For a single input we run the base and one SFT adapter, convert
each raw score to a probability with a calibrator fit only on a development split, and report the fixed
equal-weight average defined later as \Cref{eq:comp}. The $\tfrac12$ weights are fixed in advance (never
tuned on test rows), so the only cost is the second inference pass. This is the output-space composition
of \Cref{sec:bg-composition}; the logit-average and convex-weight variants were visible during
development and are reported only as non-promotable ablations.

\subsection{Domain evaluation instruments: mortgage and ExpGuard}
\label{sec:setup-domain}

Act~IV evaluates the \emph{same four checkpoints} on two complementary regulated-domain instruments.

\paragraph{Mortgage (built in depth, dual-label).} We constructed a fixed, HMDA-grounded benchmark whose
unit is one incoming request carrying two independent labels, general-safety $G$ and mortgage-policy $D$
(\Cref{sec:bg-mortgage}). Requests are grounded in de-identified, banded HMDA fact sheets --- never
verbatim records --- through an agentic pipeline that plans, grounds, generates, adversarially mutates,
and rubric-bound-judges each item, accepting it only if its assigned label matches the target, before a
provenance-tracked, family-isolated split into the frozen release (\Cref{fig:pipeline}). The benchmark
has $994$ rows over four $G\times D$ quadrants --- G0/D0 ($450$), G0/D1 ($502$), G1/D1 ($42$), and G1/D0
($0$, empty by construction) --- with the load-bearing \textbf{G0/D1} stratum (reads safe, is a
violation) the largest. It spans seven content strata (fair-lending $204$, fraud $112$, UDAAP $90$,
disclosure $66$, ATR/QM $54$, privacy $18$, and benign $450$) and is split into train $604$ / dev $149$
/ public-test $146$ / private-test $95$. As a fairness-invariance gate it embeds $39$ protected-class
\emph{minimal pairs} ($78$ rows) that are identical except for one protected token; the induced
prediction gap is denoted $\Delta_{\text{context}}$. The row composition is in \Cref{tab:composition} and
the zero-shot base results (AP$\cdot$G, AP$\cdot$D, $\Delta_{\text{context}}$) in \Cref{tab:baseline}
(Act~IV).

\begin{takeaway}{A measuring stick, not a legal finding}
Mortgage labels are assigned by an \emph{LLM judge} against written policy cards --- \emph{not}
SME-adjudicated (self-consistency only; no human Fleiss-$\kappa$). The G1/D0 quadrant is empty and some
strata are small. The benchmark \emph{surfaces} guard behavior; it certifies nothing about any real
lender or model, and licenses no fair-lending claim.
\end{takeaway}

\paragraph{Finance, health, law (external breadth: ExpGuard).} To test whether the
specialization/transfer pattern recurs across regulated verticals \emph{with expert labels} --- a
strictly stronger labeling tier than the mortgage LLM-judge --- we also score the same four checkpoints
on ExpGuard \citep{expguard2026}, an expert-annotated moderation set of $2{,}275$ rows spanning
\textbf{finance}, \textbf{health}, and \textbf{law}, using its \code{prompt\_label} for input-prompt
classification (matching our task). This is a single-label transfer probe, not the mortgage dual-$G\times D$
construct --- one domain built in depth plus three external verticals for breadth. We report aggregate
and per-domain AP for the four base checkpoints (\Cref{tab:expguard}, \Cref{fig:expguard-domains}). The
four-checkpoint result is complete and reproducible from committed text-free per-row scores; the tuned
comparison is future work.
